\section{Conceptual Framework}

The paper separates three concepts that are often conflated. Fuel poverty identifies present affordability hardship under current prices and required energy needs. Energy burden measures required or actual energy cost relative to residual household income. Household fuel price risk captures future exposure, sensitivity, and limited adaptive capacity under price-shock scenarios.

\begin{equation}
\text{Household fuel price risk} = \text{exposure} + \text{sensitivity} - \text{adaptive capacity}.
\end{equation}

The implemented baseline burden indicators should therefore be read as current-price affordability measures, not as the final fuel-price-risk outcome. They provide the income--cost foundation for later shock scenarios, but they do not themselves represent future exposure unless fuel-price, tariff, and intervention counterfactuals are introduced.

In the empirical work, this distinction is operationalised through two related but non-equivalent burden measures. The first is a deterministic baseline required-energy-burden ratio: SAP-derived required energy cost divided by MSOA after-housing-cost disposable income. The second is a probabilistic threshold measure: the probability that a household drawn from the property's LSOA income distribution would face a required energy burden above 10 percent. The former is intuitive as a cost-to-income ratio; the latter is better suited to identifying local income-tail exposure. Neither measure is an official fuel-poverty classification, and neither should be read as future household fuel price risk before price shocks are modelled.

This matters for interpretation. A top-decile modelled burden-risk group identifies the high end of Westminster's current required-cost and income-distribution assumptions. It does not mean that every property in that group is occupied by a verified fuel-poor household, nor that all other properties are protected. Instead, it marks where required energy demand, local income sensitivity, and dwelling characteristics make high burden more probable under current prices. The future-facing risk analysis must then ask how this baseline exposure changes when gas prices, electricity prices, standing charges, tariff access, and policy mitigations are varied.

This distinction matters because decarbonisation is not the same as de-risking. Electrification can reduce emissions before it reduces household price risk where electricity prices remain linked to gas through wholesale market design. The protective value of heat pumps, PV, storage, smart tariffs, and flexibility therefore depends on the electricity-gas price ratio, tariff access, capital access, dwelling form, tenure, and governance.
